\subsection{Saturation and the Born--Infeld candidate completion}
  \label{subsec:stepB-BI}

  The IR expansion~\eqref{eq:EH-result} is valid only for
  $\ell_c \gg \ell_\chi$.
  In the ultraviolet regime $\ell_c \sim \ell_\chi$,
  higher-derivative corrections become comparable to the
  Einstein--Hilbert term, and the expansion breaks down.
  The saturation bound $|\partial_t \chi| \leq c_{\mathrm{BI}}$
  provides the principle that constrains the possible UV
  completions.

  As formulated in~\cite{Beau2026c}, minimal conditions
  for a UV completion compatible with bounded flux propagation
  are:
  \begin{enumerate}[label=(\roman*)]
    \item it reduces to a quadratic (Einstein--Hilbert-like)
    form in the IR ($\ell_c \gg \ell_\chi$),
    \item it enforces a strict upper bound on admissible field
    invariants at $\ell_c \sim \ell_\chi$,
    \item it introduces no additional microscopic scales
    beyond $c_{\mathrm{BI}}$ and $\hbar_\chi$.
  \end{enumerate}

  Applied to the gravitational sector, the Born--Infeld
  candidate completion takes the form:
  \begin{equation}
    S_{\mathrm{coh}}^{\mathrm{full}}[g, \psi]
    = \frac{c_{\mathrm{BI}}^4}{\hbar_\chi^2}
    \int_{\mathcal{M}} \left(
                         \sqrt{-\det\!\left( g_{\mu\nu}
                                        + \frac{\hbar_\chi}{c_{\mathrm{BI}}^2} R_{\mu\nu} \right)}
                         - \sqrt{-g}
    \right) d^4x
    + S_{\mathrm{matter}}[g, \psi].
    \label{eq:BI-gravity}
  \end{equation}

  In the IR expansion $\hbar_\chi / c_{\mathrm{BI}}^2 \ll \ell_c^2$:
  \begin{align}
    \sqrt{-\det\!\left( g_{\mu\nu}
                   + \frac{\hbar_\chi}{c_{\mathrm{BI}}^2} R_{\mu\nu} \right)}
    &\approx \sqrt{-g}
    \left(
      1
      + \frac{\hbar_\chi}{2c_{\mathrm{BI}}^2} R
    + \mathcal{O}\!\left(\frac{\hbar_\chi^2}{c_{\mathrm{BI}}^4}
                     R_{\mu\nu} R^{\mu\nu}\right)
    \right),
    \label{eq:BI-expand}
  \end{align}
  so that:
  \begin{equation}
    S_{\mathrm{coh}}^{\mathrm{full}}
    \;\xrightarrow{\;\ell_c \gg \ell_\chi\;}\;
    \frac{c_{\mathrm{BI}}^2}{2\hbar_\chi}
    \int_{\mathcal{M}} R \sqrt{-g}\, d^4x
    + S_{\mathrm{matter}}
    + \mathcal{O}(\ell_c^{-4}).
    \label{eq:BI-IR}
  \end{equation}

  \begin{proposition}[Born--Infeld structure as an
  admissible gravitational UV completion]
    \label{prop:BI-unique}
    Under conditions (i)--(iii), ontological uniqueness, and
    Hypothesis~[H-ext] (admissible coherence extensivity),
    the Eddington-inspired Born--Infeld
    form~\eqref{eq:BI-gravity} with
    $\mathcal{R}_{\mu\nu} = R_{\mu\nu}$ is an admissible UV
    completion of the Einstein--Hilbert term in
    $S_{\mathrm{coh}}$; the admissible completions form an
    explicitly parametrised family, within which the
    determinantal member is not uniquely selected.
  \end{proposition}

  \begin{proof}
    The classification is established in~\cite{Beau2026h}:
    under~[H-ext] the admissible completers take the
    spectral form $\exp\sum_i h(\lambda_i)$, with the sector
    function $h$ constrained only at the origin and at the
    saturation boundary; the determinantal member
    corresponds to $h(\lambda) = \tfrac12\log(1+\lambda)$,
    and an explicit one-parameter family of distinct
    admissible completers realises the residual freedom.
    A classification of two-derivative covariant
    endomorphisms selects $g^{\mu\alpha}R_{\alpha\nu}$ as
    the unique admissible spectral carrier, and the trace
    ambiguity
    $R_{\mu\nu} \mapsto R_{\mu\nu} - \alpha\,g_{\mu\nu}R$
    is resolved by matching to the $a_2$ heat-kernel
    coefficient of~\cite{Beau2026h}, selecting
    $\alpha = 0$.
    The input a unique selection would require is stated
    in~\cite{Beau2026h}.
  \end{proof}
